\documentclass[chiro,pdf]{prosper}
\usepackage{url}

% Toegegeven, LaTeX is misschien niet de meest laagdrempelige manier
% om slides te maken, maar het laat zich wel heel gemakkelijk
% in versiecontrole gieten.


\title{GroepsAdministratieProgramma}
\author{Introductie, stand van zaken en demo}
\email{informatica@chiro.be}
\institution{
	Chirojeugd Vlaanderen \\
	Kipdorp 30 \\
	2000 Antwerpen
}

\slideCaption{\textit{Introductie, stand van zaken en demo}}
\Logo{\includegraphics[height=15mm]{images/logo.eps}}


\begin{document}

\maketitle

%
% OVERZICHT
%

\overlays{5}
{
\begin{slide}{Overzicht}

\begin{itemstep}
\item Afbakening project
\item Demonstratie beta-2
\item Wat verandert er in Kipadmin
\item Te doen 
\item Vragen
\end{itemstep}

\end{slide}
}

%
% PROJECTAFBAKENING
%

\overlays{6}
{
\begin{slide}{Projectafbakening}

De functionaliteit van het nieuwe programma wordt teruggebracht tot de essentie:
\begin{itemstep}
\item Beheer van leden voor de groep
\item Aansluiten van leden en leiding
\item Verzekering loonverlies
\item Dubbelpuntabonnementen
\item Bivakaangifte
\item E-mailadressen in en uit maillijsten
\end{itemstep}

\end{slide}
}

%
% DEMO
%

\overlays{12}
{
\begin{slide}{Demo}

\begin{itemize}
\item Webtoepassing: \url{https://gap.chiro.be/gap}
\onlySlide*{2}
{
\begin{itemize}
\item geen installatie nodig
\item meerdere personen kunnen tegelijk werken
\item platformonafhankelijk (Windows, Linux, Mac, \ldots)
\end{itemize}
}
\fromSlide{3}{\item Kijken wie er aangesloten is}
\fromSlide{4}{\item Personen wijzigen}
\fromSlide{5}{\item Personen aansluiten}
\fromSlide{6}{\item Automatische afdelingsverdeling}
\fromSlide{7}{\item Lijsten en etiketten via Excel}
\fromSlide{8}{\item Eigen categorie\"en en functies}
\fromSlide{9}{\item Adressensysteem}
\onlySlide*{10}{
\begin{itemize}
\item meerdere adressen per persoon
\item ingebouwde stratenlijst
\end{itemize}
}
\fromSlide{11}{\item Gezinnen beheren}
\onlySlide*{12}{
\begin{itemize}
\item broertjes en zusjes maken
\item een heel gezin tegelijk verhuizen
\end{itemize}
}
\end{itemize}

\end{slide}
}


%
% WAT VERANDERT ER IN KIPADMIN
%

\overlays{4}
{
\begin{slide}{Wat verandert er in Kipadmin}

\begin{itemize}
\fromSlide{2}{\item Automatische aanmaak van `aansluitingen' voor leden waarvan de probeerperiode voorbij is}
\fromSlide{3} {\item \'Alle leden in Kipadmin; niet enkel keti's en aspi's}
\fromSlide{4} {\item Bijwerken van de stratenlijst zal op het secretariaat moeten gebeuren.}
\end{itemize}

\end{slide}
}


%
% TODO
%

\overlays{7}
{
\begin{slide}{Te doen}
\begin{itemize}
\fromSlide{2}{\item Logica probeerperiode}
\fromSlide{3}{\item Nieuw werkjaar}
\fromSlide{4}{\item Verzekering loonverlies}
\fromSlide{5}
{
\item Communicatie GAP $\leftrightarrow$ Kipadmin
\begin{itemize}
\item gegevens van GAP naar Kipadmin
\item AD-nummers van Kipadmin naar GAP
\end{itemize}
}
\fromSlide{6}{\item Inschrijvingen voor bivak}
\fromSlide{7}{\item Interface om stratenlijst bij te werken}
\end{itemize}
\end{slide}
}


%
% VRAGEN
%

\begin{slide}{Vragen?}
\end{slide}

\end{document}
